% ============================================================
% PROPOSAL: OP-Q11 spin-statistics — Euclidean OS route
% Version: v2 (2026-06-10). STATUS: DRAFT FOR GPT REVIEW.
% Supersedes v1. All review-round-1 corrections incorporated.
% NOT inserted into ECT_preprint.tex.
% Companion document: OP_Q11_converse_calculation_v3.tex
% ============================================================

% ------------------------------------------------------------
% BLOCK 1 — Status rewrite (sec:fermion_statistics)
% ------------------------------------------------------------

\textit{Status update: the present subsection establishes the
structural/topological route only (exchange topology,
representation compatibility).
The Euclidean Osterwalder--Schrader route developed below upgrades
only the \emph{free spinor-kernel} statement (Q11.1a, candidate
Level~A as a mathematical statement) together with its conditional
ECT application (Q11.1b, Level~B); it does not upgrade the full
fermionic sector, which remains open (Q11.4).}

ECT identifies a native Euclidean route to the spin--statistics
connection: in a free minimal Euclidean $Spin(4)$ Dirac kernel, OS
positivity is compatible with the Grassmann/antisymmetric
quantisation and excludes the corresponding commuting/symmetric
half-integer spinor quantisation under the assumptions stated
below.
Applying this route to ECT fermions remains conditional on deriving
or verifying the ECT spinor quadratic kernel and its
reflection-positive structure.

% ------------------------------------------------------------
% BLOCK 2 — Proposition Q11.1a (test form) + assumptions +
%           graded-form remark + Q11.1b separation
% ------------------------------------------------------------

\paragraph{Proposition Q11.1a: OS spin--statistics test for a free
minimal Euclidean $Spin(4)$ Dirac kernel (candidate Level~A;
mathematical statement).}
For a local, minimal, first-order Euclidean Dirac kernel
transforming in a half-integer representation of
$Spin(4)\simeq SU(2)_L\times SU(2)_R$, with a specified OS
reflection map, reflection matrix, Euclidean adjoint, and positive
mass gap, the Grassmann Gaussian construction gives a
reflection-positive reconstruction.
Conversely, the corresponding commuting/symmetric quantisation of
the same half-integer local spinor kernel fails reflection
positivity.

Here the ``OS reflection structure'' means:
\begin{enumerate}[label=(\roman*)]
  \item Euclidean covariance of the quadratic kernel;
  \item a reflection $\Theta:w\mapsto-w$ compatible with the
    ordered-branch convention $t=w$;
  \item a specified spinorial reflection matrix implementing
    $\Theta$ on the $Spin(4)$ representation;
  \item locality/temperedness of the two-point kernel;
  \item support of test functions in the positive half-space $w>0$;
  \item the usual OS sesquilinear form built from the reflected
    adjoint.
\end{enumerate}
The statistics assignment itself is \emph{not} included among these
assumptions: the comparison is between the Grassmann/antisymmetric
Gaussian construction and the commuting/symmetric construction for
the \emph{same} quadratic spinor kernel.

For fermionic polynomials $F$ the OS form
$\langle(\Theta F)^{*}F\rangle_E$ is the \emph{graded} OS form:
$\Theta$ acts anti-linearly, reflects the support, applies the
spinorial reflection matrix, and reverses the order of Grassmann
factors with the corresponding signs.
The scalar RP formula of
\S\ref{sec:qs_reflection_positivity} is not being used unchanged.

\emph{Scope restriction (kernel class).}
For the minimal Dirac kernel considered here, the leading local
half-integer quadratic operator has the first-order
$\gamma^E\!\cdot\!\partial+m$ structure.
More general local higher-derivative spinor kernels would require a
separate RP analysis and are not covered by Q11.1a.

\emph{Proof structure.}
(\emph{Direct.})
For the Grassmann Gaussian sector with the ECT Euclidean Dirac
operator $D_E$ and two-point kernel $S_E=D_E^{-1}$, in conventions
to be fixed jointly with \S\ref{sec:dirac_eq} (note: the present
\S\ref{sec:dirac_eq} is formulated in Lorentzian signature; the
Euclidean gamma conventions are introduced in the companion
calculation draft and must be adopted into the main text together
with this block), the graded OS form for $F$ supported at $w>0$
reduces by half-space contour evaluation to a fixed-spatial-momentum
form with kernel proportional to a non-negative matrix (a projector
identity), and is therefore non-negative.
(\emph{Converse.})
For the commuting/symmetric quantisation of the same kernel, the
reflected two-point form at fixed spatial momentum reduces to a
matrix with strictly negative modes; an explicit negative-norm test
function is exhibited.
Both computations are carried out in the companion draft
\texttt{OP\_Q11\_converse\_calculation\_v3.tex} (review status:
confirmed conditionally in the trial OS convention; pending the
final convention check recorded there).
Safe status wording: the free-kernel OS spin--statistics test has
a completed candidate calculation; its application to ECT fermions
remains conditional on deriving or verifying the ECT spinor kernel
and spinorial reflection-positive OS structure.

Thus Q11.1a may become a Level~A mathematical theorem once the
Euclidean spinor/OS conventions are fixed and the convention-matching
check is completed, but Q11.1b is
not Level~A unless the ECT spinor quadratic kernel itself is
obtained from the condensate construction and shown to satisfy the
stated OS assumptions.
Its ECT application is conditional until then (Q11.1b, Level~B
conditional); the Dirac form of \S\ref{sec:dirac_eq} currently
enters with Level~B inputs, and the reflection-positivity results of
\S\ref{sec:qs_reflection_positivity} cover the free quadratic
scalar subsectors only.
The composite/effective character of the ECT fermionic fields does
not invalidate the route, but it precludes claiming a
first-principles Level~A ECT derivation of fermionic statistics at
this stage.

% ------------------------------------------------------------
% BLOCK 3 — Q11.2 descent conditions (defect clause restricted)
% ------------------------------------------------------------

\paragraph{Q11.2: descent to the ordered branch (Level~B
conditional).}
The $Spin(4)\to Spin(3)$ breaking pattern carries the half-integer
representation content into the ordered Lorentzian branch.
This group-theoretic descent is not by itself a spin--statistics
theorem.
The descent additionally requires:
\begin{enumerate}[label=(\roman*)]
  \item compatibility of the OS reflection $\Theta$ with the
    ordered-branch time $t=w$;
  \item continuity of the half-integer representation content
    across the breaking pattern $Spin(4)\to Spin(3)$;
  \item ordinary $3{+}1$-dimensional exchange topology,
    $\pi_1(\mathcal M_2)=\mathbb Z_2$, with no effective anyonic
    sector for fundamental exchange in the primary branch;
  \item for the pointlike primary-branch excitations considered in
    the spin--statistics statement, no topological-defect,
    skyrmion, or texture sector changes the exchange class;
    extended/solitonic sectors are outside Q11.2 and remain
    separately classified;
  \item inherited locality/domain-of-dependence and
    positive-Hilbert-space reconstruction in the ordered branch;
  \item the proof is performed in the symmetric Euclidean/OS sector
    and transported into the ordered branch; it is not derived from
    the broken $O(3)$ structure alone.
\end{enumerate}

% ------------------------------------------------------------
% BLOCK 3b — Q11.3
% ------------------------------------------------------------

\paragraph{Q11.3: CAR algebra and effective Fock reconstruction
(Level~B conditional).}
Given Q11.1b and Q11.2, the OS reconstruction applied to
antisymmetric Schwinger functions yields the CAR on the
reconstructed one-particle space; in an orthonormal mode basis this
becomes $\{\hat a_i,\hat a_j^{\dagger}\}=\delta_{ij}$.
The result is conditional on locality, cluster decomposition, and
positive-Hilbert-space reconstruction, and addresses
items~(ii)--(iii) of the ``what ECT does not yet derive'' list of
the fermionic statistics subsection at conditional Level~B, not as
a standalone result.

% ------------------------------------------------------------
% BLOCK 4 — Status-table rows (Statement & Status & Comment)
% ------------------------------------------------------------

Q11.1a: OS spin--statistics test for a free minimal
  Euclidean $Spin(4)$ Dirac kernel
  & A candidate (mathematical; converse calculation completed in
    trial OS convention)
  & Pending final convention check against the OS/Fermi literature
    and the Euclidean Dirac block to be adopted in the main text;
    minimal first-order kernel, $m>0$; Majorana/charge-conjugated
    reformulation would require translation \\
Q11.1b: application to the ECT spinor quadratic
  subsector
  & B conditional
  & Requires derivation or explicit verification of the ECT spinor
    kernel and spinorial reflection positivity; composite/effective
    fermions prevent Level-A ECT status at this stage \\
Q11.2: descent $Spin(4)\to Spin(3)$ into the
  ordered branch
  & B conditional
  & Pointlike primary-branch excitations only; conditions
    (i)--(vi); extended/solitonic sectors separately classified \\
Q11.3: CAR/Fock reconstruction via OS
  & B conditional
  & Conditional on Q11.1b + Q11.2 \\
Q11.4: full interacting fermionic/coherent sector
  & Open
  & Inherits the Level~B/open status of interacting reflection
    positivity \\

% ------------------------------------------------------------
% PLACEMENT PLAN (unchanged from v1; no edits performed)
% + forbidden phrasing: "ECT derives spin-statistics from Euclidean
%   ontology" must not appear.
% + NEW dependency surfaced in v2: main text currently has NO
%   Euclidean gamma conventions (sec:dirac_eq is Lorentzian);
%   adopting Q11.1a requires a Euclidean-conventions block.
% ============================================================
